\section{Vectors V4 and V5: Presentation and Version Selection}
\label{sec:v45}

\subsection{Vector V4: Selective Result Presentation}
\label{sec:v4}

\subsubsection{Attack Description}

The \bisect{} algorithm produces a single deterministic plan given a fixed
random seed (or uses \cs{} to find the optimal seed).
However, a hostile actor running the algorithm multiple times with
different seeds, or across different convergence thresholds, will
observe a distribution of outputs.
The attack consists of selecting the run with the most partisan-favorable
outcome and presenting only that run to courts and the public.

This attack does not affect the plan geometry for any individual run;
it exploits the selection process across runs.
It is analogous to publication bias in scientific research:
by reporting only the most favorable result from many trials, the actor
presents a biased sample of the algorithm's output distribution.

\subsubsection{Effect Estimate}

From B.7's seed variance analysis~\citep{b7paper}, the standard deviation
of Democratic seat outcomes across random seeds is $\sigma \approx 0.8$
seats nationally for a 50-state run.
If an adversary runs the algorithm $m$ times and selects the most favorable
outcome, the expected maximum shift is:
\[
  \Delta D_{V4}(m) \approx \sigma \cdot \Phi^{-1}\!\left(\frac{m-1}{m+1}\right),
\]
where $\Phi^{-1}$ is the standard normal quantile.
For $m = 10$ runs: $\Delta D_{V4}(10) \approx 0.8 \times 1.54 \approx 1.2$ seats.
For $m = 100$ runs: $\Delta D_{V4}(100) \approx 0.8 \times 2.33 \approx 1.9$ seats.

This is larger than V1--V3 in absolute terms but does not change any
individual plan's geometry.
Whether V4 constitutes gerrymandering depends on whether the act of selection
is itself the manipulative act.

\subsubsection{Detection Mechanism}

The \dia{} pre-registration protocol eliminates V4 by requiring that the
plan submitted for certification be the \emph{first} run under the
pre-registered seed (or the \cs{} optimal plan, which is deterministic).
The manifest's SHA-256 chain locks the submitted plan to a specific seed
and run configuration.
Any evidence that additional runs were performed and the results withheld
constitutes a protocol violation.
Detection probability: $\delta_4 = 0.80$ (limited by the ability to
audit whether additional runs were performed without access to the
computational logs).

\subsection{Vector V5: Algorithm Version Selection}
\label{sec:v5}

\subsubsection{Attack Description}

The \bisect{} binary evolves over time; new versions may produce
different outputs from the same inputs due to algorithmic improvements,
bug fixes, or interface changes.
A hostile actor could run multiple versions of the algorithm, observe
partisan outcomes for each, and deploy the version producing the most
favorable outcome while claiming the version selection was a routine
software update decision.

\subsubsection{Effect Estimate}

Version-to-version variation in the \bisect{} binary is bounded by
the B.0 bakeoff result~\citep{b0paper}: within the GeoSection algorithm
family (which is the statutory default), version changes affect only
implementation details (optimization speedups, tie-breaking rules) but
not the fundamental algorithm structure.
We estimate $\Delta D_{V5} \leq 0.3$ seats nationally for version changes
within the same algorithm family.
Cross-family changes (e.g., switching from GeoSection to ApportionRegions)
produce much larger effects (up to $\sim$18 seats) but are structurally
detectable as algorithm-class changes.

\subsubsection{Detection Mechanism}

The manifest includes a SHA-256 hash of the \bisect{} binary itself.
Any change to the binary changes the hash.
The \dia{} protocol requires that the deployed binary hash match the
hash of the statutorily approved algorithm version.
Detection probability: $\delta_5 = 1.00$ (cryptographic guarantee).
Binary version selection is thus the easiest V5 attack to prevent.
